Let's take a look at the following 2D convolution kernel:
\begin{equation}
    W(r)=k\begin{cases}
        1-6(\frac{r}{h})^2+6(\frac{r}{h})^3 
        & \textnormal{for }0\leq\frac{r}{h}<\frac{1}{2} \\
        2(1-\frac{r}{h})^3
        & \textnormal{for }\frac{1}{2}\leq\frac{r}{h}<1 \\
        0 
        & \textnormal{otherwise} \\
    \end{cases}
\end{equation}

\subsection{Calculate the normalization factor $k$ such that the 
kernel is normalized to unity.}
    Normalization condition:
    \begin{align}
        1
        &\overset{!}{=}\int_{\mathbb{R}^2}W(|\vec r|)\cdot\dd^2\vec r \\
        &=\int_0^{2\pi}d\varphi\cdot\int_0^\infty W(r)\cdot r \dd r \\
        &=2\pi k\cdot\bigg[
            \int_0^{h/2}\bigg(
                1-6\bigg(\frac{r}{h}\bigg)^2+
                6\bigg(\frac{r}{h}\bigg)^3
            \bigg)\cdot r \dd r+2\int_{h/2}^h\bigg(1-\frac{r}{h}\bigg)^3\cdot r \dd r
        \bigg] \\
        &=2\pi kh^2\cdot\bigg[
            \int_0^{1/2}\bigg(
                x-6x^3+
                6x^4
            \bigg) \dd x+2\int_{1/2}^1 x - 3x^2 + 3x^3 - x^4 \dd x
        \bigg] \ \ \ \ \ \ \ \ \textnormal{with $x:=\frac{r}{h}$} \\
        &=2\pi h^2k\cdot\bigg[
            \bigg(
                \frac{1}{2}x^2-\frac{3}{2}x^4 + \frac{6}{5}x^5 
            \bigg)\bigg|_0^{1/2}+2\bigg(
                \frac{1}{2}x^2 - x^3 + \frac{3}{4}x^4 - \frac{1}{5}x^5
            \bigg)\bigg|_{1/2}^{1}
        \bigg] \\
        &=2\pi h^2k\cdot\bigg[
            \frac{1}{8} - \frac{3}{32} + \frac{3}{80} + 2 \cdot \left(\frac{1}{2}- 1 + \frac{3}{4} -\frac{1}{5} - \frac{1}{8} + \frac{1}{8} -\frac{3}{64} + \frac{1}{160}\right)  
        \bigg] \\
        &=\frac{7}{40}\pi h^2k
    \end{align}
    Therefore, the normalization constant $k$ has to be
    \begin{equation}
        k=\frac{40}{7\pi h^2}
    \end{equation}

\newpage
\subsection{Use the supplied code to smooth an image}
    We have to modify the code at two positions.
    First, we initialize the kernel:
    \begin{lstlisting}
        # now set the values of the kernel
        for i in np.arange(pixel):
            for j in np.arange(pixel):
                # TODO: do something sensible here to set the real part of the kernel
                r = np.sqrt(i**2 + j**2)
                x = r / hsml
                if 0 <= x < .5:
                    kernel_value = 1 - 6*x**2 + 6*x**3
                elif .5 <= x < 1:
                    kernel_value = 2 * (1 - x)**3
                else:
                    kernel_value = 0
        
                kernel_real[i][j] = k * kernel_value\end{lstlisting}
    Then, at a later point in the supplied code, 
    we multiply the colors (in $k$-space) with the kernel:
    \begin{lstlisting}
        # multiply with kernel in Fourier space
        # TODO: fill in code here
        color_kspace *= kernel_real\end{lstlisting}
    The results can be seen below. On the left is the original image, 
    and on the right is the smoothed one.
    \begin{figure}[h!]
        \centering
        \begin{minipage}{.5\linewidth}
          \centering
          \subfloat[original image]{
            \label{:a}
            \includegraphics[scale=.45]{../figures/aq-original.pdf}
          }
        \end{minipage}%
        \begin{minipage}{.5\linewidth}
          \centering
          \subfloat[smoothed image (with $h=10$px)]{
            \label{:b}
            \includegraphics[scale=.45]{../figures/aq-smoothed.pdf}
          }
        \end{minipage}
        % \caption{}
        % \label{}
    \end{figure} \ \\

\newpage
\subsection{Calculate the sum of all channels}
    To check our results, we determine the sum of all red, green 
    and blue pixel values both before and after the smoothing.
    For this, we define the following function:
    \begin{lstlisting}
        def print_color_sums(img):
            r_sum, g_sum, b_sum, total_sum = 0, 0, 0, 0
            for row in img:
                for px in row:
                    r, g, b = px[0], px[1], px[2]
                    r_sum += r
                    g_sum += g
                    b_sum += b
                    total_sum += r + g + b
            print('r:', int(r_sum))
            print('g:', int(g_sum))
            print('b:', int(b_sum))
            print('total:', int(total_sum))\end{lstlisting}
    Calling this function both before the smoothing process prints out
    the following values: \\
    $$
        \textnormal{r: 25182443} \ \ \ \ \
        \textnormal{g: 12287845} \ \ \ \ \
        \textnormal{b: 24795222} \ \ \ \ \
        \textnormal{total: 62265510} \ \ \ \ \
    $$
    And afterwards:
    $$
        \textnormal{r: 8127774} \ \ \ \ \
        \textnormal{g: 3965970} \ \ \ \ \
        \textnormal{b: 8002796} \ \ \ \ \
        \textnormal{total: 20096540} \ \ \ \ \
    $$
    As can also be seen by looking directly at the images, the total sum of 
    pixel values decreases during the smoothing, i.e. the image gets darker. In total we lost a two third of the orginal, but we don't know why.
    % Thus, smoothing the image leaves both the sum of pixel values of a 
    % given color as well as the total sum of pixel values unchanged.
